\subsection{ALeRCE Aggregate Calibration}
\label{sec:results:alerce_agg}

Figure~\ref{fig:alerce_reliability} shows the reliability diagram for 
ALeRCE's light curve classifier evaluated on 1,576 spectroscopically 
confirmed objects (30 TDE excluded; see Section~\ref{sec:results:tde}). 
The classifier exhibits systematic underconfidence across all confidence 
bins: every point lies above the perfect calibration diagonal, indicating 
that predicted probabilities are consistently lower than empirical 
accuracies.

\begin{figure}
\centering
\includegraphics[width=\columnwidth]{../figures/fig1_reliability_diagram}
\caption{Reliability diagram for ALeRCE's light curve classifier 
evaluated on 1,576 spectroscopically confirmed ZTF transients 
(30 TDE excluded). Every point lies above the perfect calibration 
diagonal, indicating systematic underconfidence across all confidence 
bins. The histogram shows sample counts per bin. 
$\mathrm{ECE} = 0.259$ (95\% CI: $[0.237, 0.280]$).}
\label{fig:alerce_reliability}
\end{figure}

The aggregate ECE is $0.259$ (95\% CI: $[0.237, 0.280]$), well above 
the 0.05 threshold for practically significant miscalibration. The 
overall accuracy is $0.709$ while the mean top-1 confidence is $0.451$, 
a gap of $+0.258$ that is consistent with the ECE estimate. This gap 
is directional and consistent — not a consequence of random binning 
artefacts — indicating a systematic mechanistic suppression of 
classifier confidence.

To assess whether this miscalibration arises from probability dilution 
across ALeRCE's 15-class output space, we renormalise the probability 
vectors onto the 4 evaluated classes and recompute ECE. The renormalised 
ECE is $0.251$ (95\% CI: $[0.230, 0.272]$), a difference of $0.008$ 
from the 15-class estimate. Probability dilution accounts for only 
3\% of the observed miscalibration; the remainder is intrinsic to the 
classifier.

\subsection{ALeRCE Per-Class Calibration}
\label{sec:results:alerce_perclass}

Figure~\ref{fig:alerce_perclass} shows per-class reliability diagrams 
for the four ALeRCE-classified transient types. The results reveal a 
striking inverse relationship between classification accuracy and 
calibration error.

SNIa, the most accurately classified class (accuracy $= 0.831$), 
exhibits the worst calibration ($\mathrm{ECE} = 0.377$, 
95\% CI: $[0.352, 0.401]$). SNIbc shows the same pattern 
(accuracy $= 0.828$, $\mathrm{ECE} = 0.362$, 
95\% CI: $[0.315, 0.405]$). Both classes are classified correctly 
more than 80\% of the time, yet the classifier assigns them 
confidences far below this accuracy level.

SNII presents the inverse case: accuracy is only $0.495$, yet ECE 
is $0.103$ (95\% CI: $[0.066, 0.141]$) — the best-calibrated class. 
The classifier is appropriately uncertain about objects it 
classifies poorly. SLSN shows intermediate behaviour 
($\mathrm{ECE} = 0.245$, 95\% CI: $[0.181, 0.319]$), with wider 
confidence intervals reflecting the smaller sample size ($n = 55$).

This pattern is mechanistically coherent. ALeRCE's light curve 
classifier is based on Random Forest \citep{Sanchez-Saez2021}, 
whose probability outputs are averages of hard leaf-node votes. 
In class-imbalanced datasets, this averaging systematically 
suppresses predicted probabilities for well-represented classes 
\citep{Niculescu-Mizil2005, Zadrozny2002}. SNIa comprises 47\% 
of our sample; its dominance mechanically drives down predicted 
probabilities across the board, with the strongest suppression 
for the class the classifier discriminates most reliably.

\begin{figure}
\centering
\includegraphics[width=\columnwidth]{../figures/fig2_perclass_reliability}
\caption{Per-class reliability diagrams for ALeRCE's light curve 
classifier. SNIa and SNIbc (accuracy $> 0.83$) show severe 
underconfidence (ECE $> 0.36$). SNII (accuracy $0.495$) is the 
best-calibrated class (ECE $= 0.103$). SLSN shows wider confidence 
intervals reflecting the smaller sample size ($n = 55$).}
\label{fig:alerce_perclass}
\end{figure}

\subsection{TDE Classification Failure}
\label{sec:results:tde}

TDE achieves a classification accuracy of $0.000$ on our 30-object 
sample. This is not a calibration problem but a taxonomy gap: TDE 
is absent from ALeRCE's 15-class output taxonomy. The classifier 
cannot predict TDE by design, assigning all probability mass to 
other classes for every TDE in our sample.

This finding has direct operational significance. TDEs are a 
high-priority science case for LSST-era time-domain astronomy, 
representing unique probes of supermassive black hole demographics 
and accretion physics \citep{Gezari2021}. A production broker 
used by the community cannot identify them. We exclude TDE from 
all calibration analyses but flag this gap as requiring attention 
before LSST first light.

\subsection{Temperature Scaling}
\label{sec:results:tempscaling}

We apply temperature scaling (Section~\ref{sec:methods:tempscaling}) 
to ALeRCE's classifier outputs. The optimal temperature is 
$T = 0.359$, substantially below unity, confirming that the 
classifier requires confidence \textit{sharpening} — the opposite 
of the overconfidence correction for which temperature scaling 
is most commonly applied \citep{Guo2017}.

Figure~\ref{fig:tempscaling} shows the reliability diagram before 
and after scaling. Post-calibration ECE is $0.031$ 
(95\% CI: $[0.017, 0.047]$), an 88\% reduction from the 
pre-calibration value of $0.259$. The confidence intervals 
do not overlap, confirming a statistically significant improvement. 
Mean confidence increases from $0.451$ to $0.712$, now closely 
matching the empirical accuracy of $0.709$.

\begin{figure}
\centering
\includegraphics[width=\columnwidth]{../figures/fig4_temperature_scaling}
\caption{ALeRCE reliability diagrams before (left) and after (right) 
temperature scaling with $T = 0.359$. ECE reduces from $0.259$ to 
$0.031$ — an 88\% improvement. The residual dip near confidence 
$0.3$ in the right panel reflects SNII overcorrection 
(see Section~\ref{sec:discussion:tempscaling}).}
\label{fig:tempscaling}
\end{figure}

Table~\ref{tab:tempscaling} shows per-class ECE before and after 
scaling. Temperature scaling substantially improves SNIa 
($0.377 \to 0.114$) and SNIbc ($0.360 \to 0.108$), with moderate 
improvement for SLSN ($0.225 \to 0.140$). However, SNII 
\textit{worsens} from $0.099$ to $0.189$. 

This result is mechanistically interpretable. SNII was the only 
well-calibrated class before scaling: the classifier's uncertainty 
was already proportionate to its error rate. A global $T < 1$ 
uniformly increases confidence across all classes; for SNII 
this produces overconfidence where none previously existed. 
A single global temperature parameter cannot simultaneously 
optimise calibration for classes with heterogeneous miscalibration 
characteristics. Class-specific temperature parameters would be 
required for a complete fix.

\begin{deluxetable}{lcccc}
\tablecaption{ALeRCE per-class ECE before and after temperature 
scaling ($T = 0.359$). \label{tab:tempscaling}}
\tablehead{
  \colhead{Class} & \colhead{$N$} & 
  \colhead{ECE (before)} & \colhead{ECE (after)} & 
  \colhead{$\Delta$ECE}
}
\startdata
SNIa  & 757 & 0.377 & 0.114 & $-0.263$ \\
SNII  & 560 & 0.099 & 0.189 & $+0.090$ \\
SNIbc & 204 & 0.360 & 0.108 & $-0.252$ \\
SLSN  &  55 & 0.225 & 0.140 & $-0.085$ \\
\enddata
\tablecomments{Positive $\Delta$ECE indicates worsening after 
temperature scaling. SNII is the only class for which global 
temperature scaling is counterproductive.}
\end{deluxetable}

\subsection{Fink Binary Classifiers}
\label{sec:results:fink}

Figure~\ref{fig:fink_reliability} shows reliability diagrams for 
Fink's two binary classifiers evaluated on 1,368 objects. Both 
exhibit severe miscalibration, but with qualitatively different 
failure modes.

The Random Forest classifier (\texttt{rf\_snia\_vs\_nonia}) assigns 
near-zero SNIa probability to virtually all objects: mean score 
$0.011$ against a true SNIa fraction of $0.471$. The reliability 
curve is essentially flat at the base rate regardless of predicted 
confidence, indicating the classifier has no discriminative power 
for SNIa identification in this sample 
($\mathrm{ECE} = 0.464$, 95\% CI: $[0.437, 0.490]$).

\begin{figure}
\centering
\includegraphics[width=\columnwidth]{../figures/fig3_fink_reliability}
\caption{Reliability diagrams for Fink's two binary classifiers 
evaluated on 1,368 ZTF transients. Left: Random Forest 
(\texttt{rf\_snia\_vs\_nonia}), ECE $= 0.464$. The flat curve 
indicates no discriminative power for SNIa identification. Right: 
SuperNNova (\texttt{snn\_snia\_vs\_nonia}), ECE $= 0.304$. Weak 
discriminative signal persists but is severely suppressed.}
\label{fig:fink_reliability}
\end{figure}

SuperNNova (\texttt{snn\_snia\_vs\_nonia}) shows qualitatively 
better behaviour: mean score $0.222$ against true fraction $0.471$, 
with a reliability curve that rises weakly with confidence 
($\mathrm{ECE} = 0.304$, 95\% CI: $[0.278, 0.332]$). The 
classifier retains some discriminative signal but suppresses it 
substantially.

Temperature scaling applied to both Fink classifiers hits the 
upper optimisation bound ($T \approx 10$) in both cases, 
achieving only modest ECE reductions: 24\% for RF 
($0.464 \to 0.353$) and 25\% for SuperNNova 
($0.304 \to 0.229$). When the optimal $T$ hits the bound, it 
indicates the score distribution lacks the monotonic structure 
temperature scaling requires. For Fink's RF classifier, no 
post-hoc transformation of the output scores can recover 
well-calibrated probabilities; the discriminative failure is 
structural.

\subsection{Cross-Broker Comparison}
\label{sec:results:comparison}

Table~\ref{tab:comparison} summarises calibration results across 
all three classifiers. ALeRCE's light curve classifier, despite 
exhibiting the systematic underconfidence described above, is 
substantially better calibrated than either Fink classifier. 
Critically, ALeRCE's miscalibration is correctable: temperature 
scaling reduces ECE to $0.031$, below the 0.05 significance 
threshold. Fink's classifiers, particularly the RF classifier, 
exhibit structural miscalibration that post-hoc methods cannot 
adequately address.

\begin{deluxetable}{lcccc}
\tablecaption{Calibration comparison across brokers and classifiers.
\label{tab:comparison}}
\tablehead{
  \colhead{Classifier} & \colhead{ECE} & 
  \colhead{95\% CI} & \colhead{ECE (post-TS)} &
  \colhead{$T$}
}
\startdata
ALeRCE LC    & 0.259 & [0.237, 0.280] & 0.031 & 0.359 \\
Fink RF      & 0.464 & [0.437, 0.490] & 0.353 & $\sim$10 \\
Fink SNN     & 0.304 & [0.278, 0.332] & 0.229 & $\sim$10 \\
\enddata
\tablecomments{Post-TS: ECE after temperature scaling. 
$T \sim 10$ indicates the optimiser hit the upper bound, 
signifying structural miscalibration not addressable by 
temperature scaling.}
\end{deluxetable}

These findings establish that miscalibration is not uniform 
across production astronomical brokers. The differences are 
scientifically consequential: a broker with ECE $= 0.464$ 
cannot support reliable deferral decisions, while a broker 
with ECE $= 0.031$ post-calibration approaches the reliability 
required for automated follow-up prioritisation.